% Neuro-Mimetic Architecture: A Brain-Inspired Framework for Autonomous AI Orchestration
% Michael Clark — Are-Self Research
\documentclass[man,12pt,floatsintext]{apa7}

\usepackage[utf8]{inputenc}
\usepackage[T1]{fontenc}
\usepackage{amsmath,amssymb,amsfonts}
\usepackage{graphicx}
\usepackage{booktabs}
\usepackage{hyperref}
\usepackage{algorithm}
\usepackage{algpseudocode}
\usepackage{listings}
\usepackage{xcolor}

\definecolor{codebg}{rgb}{0.95,0.95,0.97}
\definecolor{codegreen}{rgb}{0.2,0.6,0.2}
\definecolor{codegray}{rgb}{0.5,0.5,0.5}
\definecolor{codepurple}{rgb}{0.58,0,0.82}

\lstset{
  backgroundcolor=\color{codebg},
  basicstyle=\ttfamily\small,
  breaklines=true,
  captionpos=b,
  commentstyle=\color{codegreen},
  keywordstyle=\color{codepurple},
  numberstyle=\tiny\color{codegray},
  stringstyle=\color{codegreen},
  frame=single,
  numbers=left,
}

\usepackage[style=apa,backend=biber]{biblatex}
\addbibresource{../../templates/references.bib}

\title{Neuro-Mimetic Architecture: A Brain-Inspired Framework for Autonomous AI Orchestration}
\shorttitle{NEURO-MIMETIC ARCHITECTURE}
\author{Michael Clark}
\affiliation{Independent Researcher \\ scipraxian}

\abstract{
We present Are-Self, an open-source AI orchestration framework whose software architecture
mirrors the functional organization of the human brain. Each component in the system maps
to a neuroanatomical region and performs an analogous computational role: the Central
Nervous System executes directed-graph workflows via spike trains, the Frontal Lobe
conducts multi-turn reasoning sessions with large language models, the Hippocampus
maintains vector-embedded long-term memory with cosine-similarity deduplication, and the
Hypothalamus routes inference requests through failover strategies with circuit-breaker
protection. Unlike metaphorical ``brain-inspired'' systems that borrow terminology without
structural correspondence, Are-Self enforces a strict mapping between software modules and
brain regions, yielding a system that is simultaneously teachable to novice users and
extensible for advanced workflows. We describe the architecture, the tick-cycle execution
model, the addon-based prompt assembly pipeline, and the real-time event system
(Synaptic Cleft) that eliminates polling. The system runs entirely on consumer hardware
via Ollama, requires no API keys for core operation, and is MIT licensed. We evaluate
the architecture's expressiveness by demonstrating multi-step autonomous workflows
including code deployment, game build pipelines, and multi-agent reasoning swarms.
}

\keywords{neuro-mimetic architecture, autonomous AI, local inference, swarm reasoning,
brain-inspired computing, Ollama, open-source AI}

\begin{document}
\maketitle

\section{Introduction}

The dominant paradigm for AI application frameworks treats language models as stateless
function calls: prompt in, completion out. Systems like LangChain, AutoGen, and CrewAI
provide abstractions for chaining these calls, but their architectures bear no structural
relationship to the cognitive processes they attempt to replicate. The result is frameworks
that are powerful but opaque---debugging a failing agent chain requires understanding
implementation-specific abstractions rather than intuitive mental models.

We propose a different approach: organize the entire orchestration system around the
functional architecture of the human brain. Not as metaphor, but as a strict structural
mapping where each software module corresponds to a neuroanatomical region and performs
an analogous computational role.

Are-Self implements this principle as a Django application suite where every app is a
brain region. The Central Nervous System fires spike trains through directed graphs of
neurons. The Frontal Lobe runs multi-turn reasoning loops. The Hippocampus stores
vector-embedded memories. The Hypothalamus selects and routes models. The Temporal Lobe
schedules work iterations. The Parietal Lobe executes tools. The Prefrontal Cortex
manages projects.

This paper makes three contributions:

\begin{enumerate}
  \item A formal description of the neuro-mimetic architecture and its mapping to
        neuroanatomy (Section~\ref{sec:architecture}).
  \item The tick-cycle execution model that drives autonomous operation
        (Section~\ref{sec:execution}).
  \item An evaluation of the architecture's expressiveness across real-world autonomous
        workflows (Section~\ref{sec:evaluation}).
\end{enumerate}

\subsection{Design Constraints}

Are-Self is designed for a specific user: a 10-year-old with no money, no powerful GPU,
and no computer science degree. Every architectural decision flows from this constraint.
The system must be free (MIT licensed), local (no cloud dependency), private (no data
leaves the machine), and comprehensible (the brain metaphor must aid understanding, not
obscure it).

\section{Related Work}

\subsection{Agent Orchestration Frameworks}

LangChain \parencite{} provides a chain-of-thought abstraction for LLM applications.
AutoGen \parencite{} introduces multi-agent conversation patterns. CrewAI defines
role-based agent teams. These frameworks excel at composability but lack structural
correspondence to cognitive processes.

\subsection{Brain-Inspired Computing}

Neuromorphic hardware (Intel Loihi, IBM TrueNorth) implements spiking neural networks
in silicon. These systems operate at the neuron level; Are-Self operates at the
brain-region level, mapping functional areas rather than individual neurons.

\subsection{Cognitive Architectures}

SOAR and ACT-R provide production-rule cognitive models. Are-Self differs in targeting
practical AI orchestration rather than cognitive simulation, and in using modern LLMs
as the reasoning substrate rather than symbolic rule systems.

\section{Architecture}
\label{sec:architecture}

\subsection{Brain Region Mapping}

Table~\ref{tab:regions} shows the mapping between Are-Self's Django applications and
their neuroanatomical counterparts.

% TODO: Insert table of brain regions -> Django apps -> computational roles

\subsection{Central Nervous System}

The CNS is a generic directed-graph executor. It knows nothing about reasoning, memory,
or language models---it simply fires spikes through neural pathways.

A \textbf{NeuralPathway} is a directed acyclic graph of \textbf{Neurons} connected by
\textbf{Axons}. Each Neuron wraps an \textbf{Effector}---a configured action that can
execute locally or be distributed across the Peripheral Nervous System. Axons are typed
by \textbf{AxonType}: FLOW (unconditional), SUCCESS (on success), or FAILURE (on failure).

A \textbf{SpikeTrain} represents a single traversal of a pathway. When started, the CNS
identifies root neurons (those with no incoming axons) and fires \textbf{Spikes} into them
via Celery tasks. Each spike carries a \textbf{blackboard}---a JSON dictionary that
accumulates data as the train traverses the graph. When a spike completes, the CNS examines
outgoing axons and dispatches the next wave.

\subsection{Frontal Lobe}

The Frontal Lobe runs \textbf{ReasoningSessions}---async while-true loops that call an LLM,
parse tool calls, execute them via the Parietal Lobe, and iterate until the agent calls
\texttt{mcp\_done()} or exhausts its focus budget.

Prompt assembly uses a four-phase \textbf{addon pipeline}:

\begin{enumerate}
  \item \textbf{IDENTIFY} --- renders the identity's system prompt template
  \item \textbf{CONTEXT} --- injects environment variables, focus/XP telemetry
  \item \textbf{HISTORY} --- reconstructs recent turns via the River of Six addon
  \item \textbf{TERMINAL} --- appends the task prompt and ``your move'' invitation
\end{enumerate}

\subsection{Hippocampus}

Long-term memory as 768-dimensional pgvector embeddings. The Hippocampus stores
\textbf{Engrams}---named facts with vector embeddings, tags, relevance scores, and
full provenance (which session created them, which spikes triggered them, which identity
owns them).

Deduplication uses cosine similarity: if a new engram is $\geq 90\%$ similar to an
existing one, the save is rejected with a pointer to the duplicate.

\subsection{Hypothalamus}

Model selection and routing. The Hypothalamus maintains a catalog of models, providers,
pricing, and capabilities. When the Frontal Lobe needs an LLM, it calls
\texttt{pick\_optimal\_model()}, which walks a \textbf{FailoverStrategy} tree:
preferred model $\rightarrow$ family matching (via vector similarity between identity
and model embeddings) $\rightarrow$ local Ollama fallback $\rightarrow$ strict failure.

Circuit breakers implement exponential cooldown on failed providers.

\subsection{Identity System}

Personas with leveling. An \textbf{Identity} is a blueprint; an \textbf{IdentityDisc}
is a deployed instance with XP, level, success/failure counts, and a 768-dim vector
embedding for semantic model matching. Every 100 XP grants a level. Levels unlock
additional focus budget.

\subsection{Temporal Lobe}

Scheduling via \textbf{Iterations} divided into \textbf{Shifts}. Each shift has a turn
limit and assigned participants (IdentityDiscs). The Temporal Lobe's metronome
(Celery Beat) ticks periodically, advancing shifts when turn limits are reached.

\subsection{Parietal Lobe}

Tool execution with \textbf{hallucination armor}. The Parietal Lobe validates tool calls
against registered \textbf{ToolDefinitions}, checking parameter types and required fields
before execution. Invalid calls return structured error messages rather than crashing.

\subsection{Prefrontal Cortex}

Project management: Epics $\rightarrow$ Stories $\rightarrow$ Tasks, with status
workflows (NEEDS\_REFINEMENT through DONE), ownership tracking, and comment threads.

\subsection{Synaptic Cleft}

Real-time event bus via Django Channels. Five neurotransmitter types---Dopamine (success),
Cortisol (error), Acetylcholine (data sync), Glutamate (streaming), Norepinephrine
(monitoring)---broadcast events to the React frontend via WebSocket. The frontend's
\texttt{useDendrite()} hook subscribes to specific receptor classes, triggering refetches
when signals fire. This eliminates all polling.

\subsection{Peripheral Nervous System}

Fleet management. NerveTerminal agents register with the system, report telemetry
(CPU, memory, storage), and receive distributed spike executions. Four distribution modes:
local, all agents, one available agent, or specific targets.

\section{The Tick Cycle}
\label{sec:execution}

Everything in Are-Self serves one loop:

\begin{enumerate}
  \item PNS (Celery Beat) ticks every $N$ seconds
  \item Temporal Lobe wakes active iteration, selects current shift
  \item CNS builds SpikeTrain, fires through neural pathway
  \item Spikes cascade through neurons, executing effectors
  \item Frontal Lobe starts reasoning session (if invoked by Frontal Lobe effector)
  \item Parietal Lobe executes tool calls; Hippocampus stores/retrieves memories;
        Hypothalamus selects models
  \item Session concludes; neurotransmitter fires
  \item Next shift tick
\end{enumerate}

This cycle enables fully autonomous operation: the system reasons, uses tools, forms
memories, and manages its own work without human intervention, bounded only by the
Focus Economy's budget constraints.

\section{The Focus Economy}

A critical design element that prevents unbounded execution. Each reasoning session begins
with a focus budget ($10 + 0.5 \times (\text{level} - 1)$). Every turn and tool call
decrements focus. Novel memory formation (engram saves with novelty $> 10\%$) grants
focus back. When focus reaches zero, the session must conclude.

This creates an incentive structure: the agent is rewarded for \emph{learning} (saving
novel memories) and penalized for verbosity (consuming turns without learning). An
additional efficiency bonus grants +1 focus and 5 XP when output stays under
$\text{level} \times 1000$ characters.

\section{Evaluation}
\label{sec:evaluation}

We evaluate the architecture along three dimensions: expressiveness (can it represent
real-world workflows?), comprehensibility (does the brain metaphor aid understanding?),
and performance (does it run on consumer hardware?).

\subsection{Expressiveness}

% TODO: Case studies of real pathways
% - Pre-Release build pipeline (5-node, multi-step)
% - Agent scanning and deployment
% - Temporal Lobe scheduled reasoning iterations

\subsection{Comprehensibility}

% TODO: User study or qualitative analysis
% - Teaching the system to non-programmers
% - Time to first successful pathway

\subsection{Performance}

% TODO: Benchmarks on consumer hardware
% - Spike train execution times
% - Reasoning session latency with 7B models via Ollama
% - Memory overhead of vector operations

\section{Discussion}

\subsection{Why Not Just a Metaphor?}

Many systems use brain terminology loosely. Are-Self's strict mapping enforces a
discipline: if a new feature doesn't map to a brain region, it likely belongs in an
existing one. This constraint prevents architectural drift and maintains the system's
teachability.

\subsection{Limitations}

The architecture currently supports one active iteration per environment. The Peripheral
Nervous System lacks sophisticated health monitoring. The Hippocampus uses flat M2M
relationships rather than the hypergraph structure proposed by Frerichs for richer
memory interconnection.

\subsection{Future Work}

PNS expansion to multiple Ollama endpoints with live agent telemetry. Hippocampus
migration to hypergraph engrams. Image and audio generation as first-class effectors.
MCP (Model Context Protocol) server and client integration for cross-system tool sharing.

\section{Conclusion}

Are-Self demonstrates that organizing an AI orchestration framework around the functional
architecture of the human brain yields a system that is simultaneously intuitive for
beginners and powerful for advanced users. The strict mapping between software modules and
brain regions provides a natural vocabulary for discussing system behavior, a principled
guide for architectural decisions, and an extensible foundation for future capabilities.

The system runs entirely on consumer hardware, requires no API keys for core operation,
and is MIT licensed---making autonomous AI accessible to anyone with a laptop and
curiosity.

\printbibliography

\end{document}
